Un dipol elèctric és un sistema de dues càrregues puntuals d'igual magnitud $Q$ i signe oposat.

\begin{apartat}{1,25}
Representeu dins del requadre adjunt les línies de camp elèctric creades per un dipol elèctric. Representeu la projecció de les superfícies equipotencials en el pla de la figura.

Orientació: per al camp elèctric dibuixeu 12 línies de camp i 3 línies equipotencials per cada càrrega.

\figura[0.55]{fig1.png}
\end{apartat}

\begin{apartat}{1,25}
El valor de la càrrega és $|Q| = 1,50\ \mu\mathrm{C}$, la càrrega positiva està situada a $-5\vec{\imath}\ \mathrm{cm}$ i la càrrega negativa està situada a $5\vec{\imath}\ \mathrm{cm}$. Calculeu el camp elèctric creat pel dipol elèctric a l'origen de coordenades i també el valor del potencial elèctric a l'origen de coordenades. Per a les magnituds vectorials podeu donar les components o el mòdul, la direcció i el sentit.
\end{apartat}

\dades{%
  $k = \dfrac{1}{4\pi\varepsilon_0} = 8,99\times 10^{9}\ \mathrm{N\,m^2\,C^{-2}}$.}
